% Generated by Listas/make_statements.py from pratica-lecture-01.md, item 8. Do not edit.
\textbf{Source.} Muñoz-Garcia \& Toro-Gonzalez (2019), \emph{Strategy and Game Theory: Practice Exercises with Answers}, Ch. 1, Exercise 9 ``Applying IDSDS in the Battle of the Sexes Game'', p.~17 {[}PDF 30{]}. Reproduced verbatim (statement only; the book's worked answer follows it).



\begin{quote}

A husband and a wife are leaving work, and do not remember which event they are attending to tonight. Both of them, however, remember that last night's argument was about attending either the football game (the most preferred event for the husband) or the opera (the most preferred event for the wife). To make matters worse, their cell phones broke, so they cannot call each other to confirm which event they are attending to. As a consequence, they must simultaneously and independently decide whether to attend to the football game or the opera.



{\def\LTcaptype{none} % do not increment counter

\begin{longtable}[]{@{}lll@{}}

\toprule\noalign{}

\emph{Husband} ~\emph{Wife} & \emph{Football} & \emph{Opera} \\

\midrule\noalign{}

\endhead

\bottomrule\noalign{}

\endlastfoot

\emph{Football} & 3, 1 & 0, 0 \\

\emph{Opera} & 0, 0 & 1, 3 \\

\end{longtable}

}



\emph{Fig. 1.22} Normal-form representation of the Battle of the Sexes game



The payoff matrix in Fig. 1.22, describes the preference of the husband (wife) for the football game (opera, respectively). Payoffs also indicate that both players prefer to be together rather than being alone (even if they are alone at their most preferred event). Apply IDSDS. What is the most precise prediction about how this game will be played when you apply IDSDS?

\end{quote}

